% !PRINTED 346
\weberpage{346}

%<<CONT
und dies reduziert sich in unserem Falle des quadratischen
Körpers darauf, daß
\weq{2}{\theta^p \equiv \theta \ (\mathrm{mod}\ \mathfrak{p})}
sein muß.

Nehmen wir $\Delta = -4$, so ist $\theta = i = \sqrt{-1}$ und die Kongruenz
 (2) wird
\[
i^{p-1} \equiv 1 \ (\mathrm{mod}\ \mathfrak{p}).
\]

Da nun $(1+i)(1-i) = 2$, also weder 2 noch $(1 \pm i)$ für
ein ungerades $p$ durch $\mathfrak{p}$ teilbar sein kann, so muß $i^{p-1} = 1$,
also $p - 1$ durch 4 teilbar sein.

Demnach folgt aus \S\,92, 2.:

1) \gesperrt{Die Kongruenz}
\[
x^2 \equiv -4 \ (\mathrm{mod}\ 4p),
\]
\gesperrt{oder, was dasselbe ist,}
\[
x^2 \equiv -1 \ (\mathrm{mod}\ p)
\]
\gesperrt{hat dann und nur dann Lösungen, wenn $p \equiv 1 \ (\mathrm{mod}\ 4)$ ist.}

Nehmen wir zweitens $\Delta = 8$, so ist $\theta = \sqrt{2}$, oder, wenn
$r = \sqrt{i}$ eine 8te Einheitswurzel ist,
\[
\theta = r + r^{-1},
\]
und die Bedingung (2) läßt sich nach dem binomischen Satze
so darstellen:
\[
r^p + r^{-p} \equiv r + r^{-1} \ (\mathrm{mod}\ \mathfrak{p}).
\]

Dies fordert:
\[
(r^p - r)(r^p - r^{-1})\, r^{-p} \equiv 0 \ (\mathrm{mod}\ \mathfrak{p}),
\]
also:
\[
r^p \equiv r \ \text{oder} \ r^p \equiv r^{-1} \ (\mathrm{mod}\ \mathfrak{p}),
\]
\[
r^{p \pm 1} \equiv 1, \quad i^{\frac{p \pm 1}{2}} - 1 \equiv 0 \ (\mathrm{mod}\ \mathfrak{p}),
\]
und daraus folgt wie oben für ein ungerades $p$:
\[
p \equiv \pm 1 \ (\mathrm{mod}\ 8).
\]

2) \gesperrt{Die Kongruenz}
\[
x^2 \equiv 2 \ (\mathrm{mod}\ p)
\]
\gesperrt{ist also dann und nur dann lösbar, wenn $p \equiv 1$ oder
$p \equiv -1 \ (\mathrm{mod}\ 8)$ ist.}

Endlich setzen wir, wenn $q$ eine ungerade Primzahl ist,
\weq{3}{\Delta = \pm q}
%CONT>>
